\section{Background}
\label{sec:background}

This section introduces the concepts needed to follow the thesis. It is deliberately
scoped to what the later experiments rely on: runway detection as a pose task,
the YOLO family, the vocabulary of the Sim2Real gap and domain adaptation, and the
LARD datasets.

\subsection{Vision-Based Runway Detection as Pose Estimation}
\label{subsec:bg-runway}
A runway projected into an approach image is well approximated by a planar
quadrilateral. Detecting it can therefore be cast either as (i) bounding-box
detection of the runway region, or (ii) localization of the four runway corners as
\emph{keypoints}. This thesis adopts the keypoint (pose) formulation: the model
predicts a runway instance together with its four ordered corner points
(e.g.\ near-left, near-right, far-left, far-right). The four corners are a compact,
geometry-preserving representation that downstream modules can turn into a relative
pose estimate through a perspective-n-point solve. A key practical consequence,
relevant throughout the thesis, is that any label-altering transformation of the
image (geometric augmentation, generative restyling) must keep the four corner
annotations geometrically consistent.

\subsection{The YOLO Family and YOLOv8-Pose}
\label{subsec:bg-yolo}
\ac{YOLO} (You Only Look Once) is a family of single-stage object detectors that
predict objects densely in a single forward pass, trading a small accuracy margin
for high inference speed suitable for onboard, real-time use. YOLOv8 extends the
family with an anchor-free, decoupled detection head and a dedicated
\emph{pose} variant that adds a keypoint-regression branch, predicting a fixed number
of keypoints per detected instance. For runway detection this branch is configured
for four keypoints. All experiments in this thesis \emph{fine-tune} a
pretrained YOLOv8-pose checkpoint rather than training from scratch, so that the
comparisons isolate the effect of the training \emph{data} and adaptation
technique rather than of the backbone initialization.

\subsection{The Sim2Real Gap and Domain Shift}
\label{subsec:bg-sim2real}
When a model is trained on a source distribution $\mathcal{D}_S$ (synthetic
imagery) and deployed on a target distribution $\mathcal{D}_T$ (real imagery),
differences between the two distributions cause a drop in performance known as the
\emph{Sim2Real gap}. It manifests as \emph{covariate shift} in low-level statistics
(texture, color, sensor noise, compression), in mid-level appearance (lighting,
shadows, atmospheric scattering) and in scene composition (which locations,
viewpoints and object scales are represented). Adverse weather compounds the shift:
rain streaks, water refraction on the lens, sun glare, fog/snow attenuation and
night illumination are rarely modeled faithfully by nominal renderers and are
under-represented in nominal synthetic data.

\subsection{Families of Domain-Adaptation Techniques}
\label{subsec:bg-da}
The techniques studied in this thesis can be placed on a spectrum according to
\emph{where} they intervene on the shift:
\begin{itemize}
    \item \textbf{Input-space augmentation.} Cheap, label-preserving transformations
    applied to the source images to synthesize target-like nuisances (e.g.\ weather).
    They do not require target data but only approximate real degradations.
    \item \textbf{Self-training / pseudo-labeling.} A form of unsupervised domain
    adaptation (\ac{UDA}) that uses model predictions on unlabeled target data as
    training targets, optionally filtered by confidence and refined by temporal
    consistency. It uses real data directly but is sensitive to confirmation bias.
    \item \textbf{Generative style transfer.} Learns or applies a mapping that
    imports target appearance onto the source images while preserving content, so
    that labels remain valid. GAN-based mappings can distort geometry; diffusion-based
    conditional transfer aims to keep structure intact.
\end{itemize}
\Cref{sec:gap,sec:weather-aug,sec:pseudo-labels,sec:style-transfer} explore these
families in this order, from the least to the most direct intervention on the gap.

\subsection{The LARD Datasets}
\label{subsec:bg-lard}
The \ac{LARD} dataset family targets runway detection on approach. The first
version (LARD-V1) provides, among others, a set of \emph{real} approach images
annotated with runway corners, which this thesis uses as the held-out real test set.
The second version (LARD-V2) refines the operational design domain (a narrower
approach cone, multi-runway airports) and, most importantly here, provides multiple
\emph{synthetic} image sources for a shared set of runway locations, generated through
different rendering and imagery pipelines. This thesis uses five such sources ---
\texttt{flsim} (Flight Simulator) and \texttt{xplane} (X-Plane), which are full 3D
simulators; \texttt{arcgis} and \texttt{bingmaps}, which are satellite-imagery based;
and \texttt{ges} (Google Earth Studio), which blends satellite imagery with 3D
rendering --- plus a \emph{combined} source that mixes all five. Because the sources
share \texttt{location\_id}s, images from different sources correspond to the exact
same runway views, which is what makes the cross-source comparison in \Cref{sec:gap} a
controlled experiment. The precise composition and the sampling scheme used to keep the
training size fixed are detailed in \Cref{sec:setup}.
